\section{Experiments}
\label{sec:experiments}

\subsection{Experimental Setup}
\label{sec:exp_setup}

\paragraph{Benchmark Dataset.} We evaluate on the complete \textbf{OMTG Bench}~\cite{omtg2026}, the established benchmark for One-to-Many Video Temporal Grounding. The benchmark consists of 320 manually verified query-video pairs spanning 287 untrimmed videos from Charades, ActivityNet, QVHighlights, VTimeLLM, and Moment10M. Durations range from $12.0\text{s}$ to $506.0\text{s}$ (mean: $161.4\text{s}$), with ground-truth instance counts ranging from 2 to 20 per query (mean: 3.67 instances).

\paragraph{Evaluation Metrics.} Following the official OMTG evaluation protocol~\cite{omtg2026}, we report:
\begin{enumerate}
    \item \textbf{Count Accuracy (C-Acc, \%) $\uparrow$}: Percentage of samples where the predicted segment count exactly matches ground-truth cardinality ($M=K$).
    \item \textbf{Temporal F1 (tF1@$\xi$, \%) $\uparrow$}: Harmonic mean of precision and recall under Hungarian bipartite matching at IoU thresholds $\xi \in \{0.3, 0.5, 0.7\}$, with unmatched predictions penalized as false alarms.
    \item \textbf{Temporal IoU (tIoU, \%) $\uparrow$}: Mean temporal intersection-over-union across matched ground-truth and predicted intervals.
    \item \textbf{Effective Temporal F1 (EtF1, \%) $\uparrow$}: The primary benchmark metric, which conditions temporal F1 on exact cardinality match ($\text{EtF1} = \text{tF1} \times \mathbf{1}(M=K)$).
    \item \textbf{Cardinality Error $\downarrow$}: Mean absolute error between predicted and ground-truth event counts ($|M - K|$).
    \item \textbf{Processed Frames and Compute Saved (\%) $\uparrow$}: Aggregate video frames processed by the 8B Vision Transformer across all 320 evaluations ($N_{\text{frames}} = \sum_i F_i$), and relative frame reduction $\text{Saved} = (1 - N_{\text{frames}}^{\text{method}} / N_{\text{frames}}^{\text{native}}) \times 100\%$.
\end{enumerate}

\paragraph{Implementation Details.} Our primary grounder is frozen TimeLens-8B~\cite{timelens2024} (Qwen3-VL-8B-Instruct backbone). The temporal scout is frozen Llama-Nemotron-Embed-VL-1B (1.04B parameters) sampled at a coarse 1.0~FPS. Hyperparameters are fixed across all benchmark videos without dataset-specific tuning: detection thresholds $\tau_{\text{high}}=0.80, \tau_{\text{low}}=0.25$ ($\S$\ref{sec:stage2}), energy baseline offset $\tau_{\text{base}}=0.30$ (Eq.~\ref{eq:excess_energy}), adaptive window coefficients $\alpha_{\text{margin}}=3.5\text{s}, \beta_{\text{gap}}=3.5\text{s}$ (Eq.~\ref{eq:adaptive_geometry}), and merge threshold $\rho=0.70$ (Eq.~\ref{eq:soft_merge}).

% -------------------------------------------------------------------------

\subsection{Grounding Performance and Computational Efficiency}
\label{sec:exp_main}

\begin{table}[t]
\centering
\caption{\textbf{Master performance and efficiency comparison on OMTG Bench ($N=320$).} All models operate zero-shot without fine-tuning. Frames denotes aggregate 8B ViT frames evaluated across the benchmark. Rows with 128 frames share identical structured prompt schemas for strictly controlled comparison. Best results are in \textbf{bold}; second-best are \underline{underlined}.}
\label{tab:main_results}
\vspace{2pt}
\resizebox{\textwidth}{!}{
\begin{tabular}{l c c c c c c c}
\hline
\textbf{Configuration} & \textbf{Frames} & \textbf{Saved$\uparrow$} & \textbf{C-Acc$\uparrow$} & \textbf{tF1@.5} & \textbf{tIoU$\uparrow$} & \textbf{EtF1$\uparrow$} & \textbf{Err$\downarrow$} \\
\hline
TimeLens-8B (Paper Prompt)~\cite{omtg2026} & 103.3k & Baseline & 0.00\% & 32.76\% & 32.38\% & 0.00\% & 2.15 \\
TimeLens-8B (JSON Prompt) & 103.3k & Baseline & 19.69\% & \textbf{47.98\%} & 46.57\% & 14.23\% & 2.07 \\
\hline
+ Naive Uniform Downsample (128f) & 41.0k & -60.3\% & 16.88\% & 45.30\% & 44.23\% & 12.29\% & 1.96 \\
+ Adaptive SGDE-64 (64f) & \textbf{20.5k} & \textbf{-80.2\%} & 27.19\% & 37.41\% & 41.22\% & 14.44\% & 2.09 \\
+ Adaptive SGDE-128 (128f) [Proposed] & 41.0k & \textbf{-60.3\%} & \underline{33.12\%} & 41.82\% & 44.76\% & \underline{19.14\%} & \textbf{1.75} \\
+ Adaptive SGDE-256 (256f) & 81.9k & -20.7\% & 31.88\% & 41.82\% & \underline{45.42\%} & \textbf{19.86\%} & \underline{1.81} \\
\hline
\emph{Upper Bound (Ground-Truth Window, 128f)} & 41.0k & -60.3\% & \textbf{33.44\%} & \underline{47.59\%} & \textbf{50.54\%} & \textbf{20.99\%} & 2.08 \\
\hline
\end{tabular}
}
\end{table}

\begin{figure}[t]
    \centering
    \includegraphics[width=0.66\textwidth]{omtg_video_length_and_pruning.png}
    \caption{\textbf{Original video duration vs. retained candidate window duration on OMTG Bench ($N=320$).} The diagonal dashed line ($W_{\mathrm{len}} = T$) denotes the unpruned baseline. For short videos ($T \le 45\text{s}$), Adaptive SGDE retains the full video span ($0\%$ pruned, red points). For longer videos ($T > 45\text{s}$), candidate window extraction removes an average of $58.1\%$ of background frames (green points, shaded area).}
    \label{fig:video_length_pruning}
\end{figure}

Table~\ref{tab:main_results} and Fig.~\ref{fig:video_length_pruning} present grounding accuracy and compute efficiency on OMTG Bench. Four primary insights emerge:

\paragraph{1. Prompt Syntax vs. Grounding Capacity.} Under the unstructured prompt used in the official OMTG study~\cite{omtg2026}, TimeLens-8B outputs conversational fragments or single moments ($0.00\%$ C-Acc and $0.00\%$ EtF1). Enforcing schema-delimited JSON syntax restores multi-span output, reaching $19.69\%$ C-Acc and $14.23\%$ EtF1. This indicates that baseline failure in prior benchmarks was partly caused by prompt formatting rather than model capacity.

\paragraph{2. Video Length Scaling and Background Removal.} As shown in Fig.~\ref{fig:video_length_pruning}, benchmark videos range from $12.0\text{s}$ to $506.0\text{s}$. On long unpruned videos, models often make false positive detections across uninformative background segments. Adaptive SGDE preserves short videos ($0\%$ pruned for $T \le 45\text{s}$) while removing an average of \textbf{58.1\% of background frames} on longer videos (rising to $69.5\%$ on videos $>180\text{s}$ and up to $85.2\%$). Removing these background segments enables Adaptive SGDE-128 to increase Count Accuracy to \textbf{33.12\%} ($+13.43\%$ over the prompt baseline) and reduce Cardinality Error from $2.07 \rightarrow \textbf{1.75}$.

\paragraph{3. Comparison with Ground-Truth Window Upper Bound.} To quantify the upper bound of the frozen TimeLens-8B grounder under a 128-frame budget, we evaluate an upper bound using a ground-truth temporal window ($100\%$ GT span retention, $0\%$ background false alarms). As reported in Table~\ref{tab:main_results}, this upper bound achieves $33.44\%$ C-Acc and $20.99\%$ EtF1. Operating completely training-free without ground-truth annotations, our proposed Adaptive SGDE-128 closely approaches this upper bound, achieving \textbf{33.12\% vs. 33.44\% C-Acc} (a $0.32\%$ difference) and \textbf{19.14\% vs. 20.99\% EtF1}, demonstrating that training-free semantic scouting effectively isolates informative event regions.

\paragraph{4. Precision-Recall Dynamics and Compute Savings.} The prompt baseline exhibits higher precision ($60.70\%$) but low recall ($40.84\%$) because it under-predicts ($0\%$ multi-event C-Acc). Adaptive SGDE exposes candidate action intervals across the focused window, increasing Recall to $44.11\%$ and reaching \textbf{19.86\% EtF1} with 256 frames ($+5.63\%$ over prompt-only baseline and $+7.57\%$ over naive downsampling). In addition, focusing 8B ViT decoding on candidate windows reduces frame evaluations by \textbf{60.3\%} (128 fixed frames vs. $322.8\text{f}$ average on raw video), offloading initial temporal search to the lightweight 1.04B scout.

% -------------------------------------------------------------------------

\subsection{Algorithmic Analysis: Window Planning and Windowing Topologies}
\label{sec:exp_ablations}

\begin{table}[t]
\centering
\caption{\textbf{Ablation of temporal window planning and windowing strategies on OMTG Bench ($N=320$, 128-frame budget).} Isolating semantic scouting, single-window bounding, and soft continuity merging.}
\label{tab:ablation_windowing}
\vspace{2pt}
\resizebox{\textwidth}{!}{
\begin{tabular}{l l c c c c c}
\hline
\textbf{Window Strategy} & \textbf{Planning Mode} & \textbf{C-Acc$\uparrow$} & \textbf{tF1@.5} & \textbf{tIoU$\uparrow$} & \textbf{EtF1$\uparrow$} & \textbf{Err$\downarrow$} \\
\hline
(a) Random Temporal Crop & Random window ($40\%$ dur) & 18.44\% & 24.29\% & 24.18\% & 3.61\% & 2.02 \\
(b) Single-Window Baseline & Single bounding window & 28.44\% & 38.26\% & 36.81\% & 16.86\% & 1.94 \\
(c) Multi-Window (No Merge) & Split on gaps ($\rho=1.0$) & 28.44\% & 38.49\% & 37.12\% & 16.81\% & 1.94 \\
(d) \textbf{Adaptive SGDE-128 (Full)} & \textbf{Adaptive window + merge ($\rho=0.7$)} & \textbf{33.12\%} & \textbf{41.82\%} & \textbf{44.76\%} & \textbf{19.14\%} & \textbf{1.75} \\
\hline
\end{tabular}
}
\end{table}

\begin{figure}[t]
    \centering
    \includegraphics[width=0.86\textwidth]{scoring_method_ablation.png}
    \caption{\textbf{Evaluation of timeline scoring and candidate extraction methods on OMTG Bench.} (a) Action span retrieval: $Z$-score normalization and energy scoring increase Ground-Truth coverage to $80.3\%$ and cut missed spans by $31.1\%$ relative. (b) Proposal continuity and overlap: achieving $95.9\%$ 1-window continuity and $48.6\%$ proposal IoU.}
    \label{fig:scoring_ablation}
\end{figure}

\paragraph{1. Semantic Guidance vs. Random Windowing.} As shown in Table~\ref{tab:ablation_windowing}, evaluating a random temporal window matching the average candidate duration ratio ($40\%$) drops EtF1 to \textbf{3.61\%} and tIoU to \textbf{24.18\%}. This shows that temporal windowing without semantic guidance degrades grounding performance.

\paragraph{2. Single-Window vs. Multi-Window Topology.} Moving from raw video to a single continuous window (strategy b) improves Count Accuracy from $16.88\% \rightarrow \textbf{28.44\%}$ and EtF1 from $12.29\% \rightarrow \textbf{16.86\%}$. However, forcing distant multi-event clusters into one window retains background frames between events, limiting tIoU to $36.81\%$. While splitting intervals without soft merging (strategy c) can fragment adjacent actions ($16.81\%$ EtF1), our proposed **Adaptive SGDE** with $70\%$ soft continuity merging (strategy d) groups proximate clusters while isolating distant intervals, improving tIoU to \textbf{44.76\%}, Count Accuracy to \textbf{33.12\%}, and EtF1 to \textbf{19.14\%}.

\begin{figure}[t]
    \centering
    \includegraphics[width=0.88\textwidth]{soft_merge_ratio_ablation.png}
    \caption{\textbf{Merge threshold ($\rho$) sweep on OMTG Bench.} (a) Continuity vs. GT coverage: setting $\rho=70\%$ retains $95.0\%$ of videos within a single continuous prompt while covering $79.71\%$ of GT action spans. (b) Background duration pruned: achieving $37.17\%$ background pruning across the benchmark (up to $85.2\%$ on videos $>180\text{s}$).}
    \label{fig:soft_merge_ablation}
\end{figure}

\paragraph{3. Timeline Scoring Function Evaluation.} As shown in Fig.~\ref{fig:scoring_ablation}, raw cosine similarity with a fixed cutoff is sensitive to baseline shifts across videos, resulting in $25.1\%$ missed action spans. Video-adaptive $Z$-score normalization with dual-threshold search ($\tau_{\text{high}}=0.80, \tau_{\text{low}}=0.25$) and energy scoring $\bar{J}(a,b)$ improves ground-truth coverage from $70.9\% \rightarrow \textbf{80.3\%}$ ($+9.4\%$), improves proposal IoU from $34.2\% \rightarrow \textbf{48.6\%}$, and reduces missed action spans by $31.1\%$ relative.

\paragraph{4. Soft Merge Span Ratio Optimization ($\rho$).} Fig.~\ref{fig:soft_merge_ablation} shows the trade-off across merge ratio thresholds $\rho \in [30\%, 100\%]$. Low thresholds ($\rho \le 40\%$) preserve all videos in 1-window prompts but prune fewer background frames ($33.3\%$). High thresholds ($\rho \ge 90\%$) or gap-only merging split $7.5\%-14.7\%$ of videos across multiple disjoint prompts, disrupting multi-event context. Setting $\rho = 70\%$ achieves a balanced trade-off: keeping \textbf{95.0\% of videos in a single continuous prompt} while capturing \textbf{79.71\% of GT action spans} and removing \textbf{37.17\% of background frames} overall (up to $85.2\%$ on videos $>180\text{s}$).
